%% Chapter 5 — Conclusion and Work Plan
\label{ch:conclusion}

\section{Current Progress}
The EMS Portal is a digital marketplace that connects clients with engineering
service providers through an intelligent multi-criteria matching mechanism. The
project is built in six phases. The first two are now complete, and the
remaining phases are designed and ordered.

\textbf{Authentication and authorization are solid.} The platform has a complete
authentication system (registration, login, logout, session refresh) backed by
short-lived JWT access tokens and hashed, family-rotated refresh tokens. The
authorization layer enforces role-based access control across three levels
(role, permission, ownership), with a shared, single-source-of-truth permission
catalogue.

\textbf{Core provider-side features are built.} Service providers can create and
manage profiles with an explicit draft, active, and suspended lifecycle. They
attach skills from a seeded, public skill catalogue. The frontend provides
complete UI for authentication and profile management.

\textbf{Infrastructure and quality practices are in place.} The system runs on a
modular, layered architecture with a migrations-first database schema (three
reviewed migrations), an enforced testing pyramid with coverage thresholds, and a
six-workflow continuous-integration pipeline. The multi-criteria matching
algorithm—the project's central research contribution—is formally specified
(Chapter~\ref{ch:methodology}) and its schema is designed, ready to
implement.

\section{Work Plan}
The remaining work takes the system from its current foundation, through the
matching engine, to a complete and evaluated platform. Table~\ref{tab:workplan}
lists the remaining phases, their main deliverables, and their dependencies.

\begin{table}[ht]
	\centering
	\caption{Work plan for the remaining thesis phases.}
	\label{tab:workplan}
	\resizebox{\textwidth}{!}{%
		\begin{tabular}{cll}
			\hline
			\textbf{Phase} & \textbf{Scope} & \textbf{Principal deliverables} \\ \hline
			3 & Service requests &
				\begin{tabular}[t]{@{}l@{}}
					Client-posted requests; status state machine \\
					(open $\to$ matched $\to$ in\_progress $\to$ completed $|$ cancelled $|$ expired)
				\end{tabular} \\ \hline
			4 & Matching engine &
				\begin{tabular}[t]{@{}l@{}}
					Implement Eq.~\ref{eq:score}; seed \texttt{basic-v1} and \texttt{intelligent-v1}; \\
					per-criterion scoring; persist ranked \texttt{match\_results} audit log
				\end{tabular} \\ \hline
			5 & Engagements \& ratings &
				\begin{tabular}[t]{@{}l@{}}
					Engagement state machine; bidirectional ratings; \\
					denormalized provider rating aggregates
				\end{tabular} \\ \hline
			6 & Administration \& review &
				\begin{tabular}[t]{@{}l@{}}
					User management; weight configuration UI; \\
					basic-vs-intelligent effectiveness comparison
				\end{tabular} \\ \hline
			Final & Evaluation &
				\begin{tabular}[t]{@{}l@{}}
					Execute the penetration-testing plan; k6 load tests \\
					for latency/response time; matching-effectiveness study
				\end{tabular} \\ \hline
		\end{tabular}%
	}
\end{table}

The order is driven by dependencies. Service requests (Phase~3) supply the
demand side that the matching engine (Phase~4) uses. The matching results feed
the engagements and ratings of Phase~5. Those ratings feed back into the matching
score's rating factor, closing the feedback loop. The admin tools of Phase~6
enable the side-by-side comparison of the \texttt{basic-v1} and
\texttt{intelligent-v1} configurations that the evaluation needs. The final
evaluation then tests the complete system along the three axes from the
objectives: functional correctness, security, and network performance.

\section{Current Limitations}
Several limitations follow from the current state of the project. The matching
engine, the central research contribution, is not yet built, so the full
request-to-match-to-engagement-to-rating flow cannot be demonstrated end to end.
There is no administrator interface: the permission model recognizes
administrators, but the dedicated screens are planned for Phase~6. Geographic
proximity is approximated using the haversine formula on plain latitude and
longitude coordinates; a dedicated geographic extension such as PostGIS is
deliberately deferred to Phase~4 to keep the early phases simple. The user
interface is currently a rough, functional MVP, with visual design left for later
work.

\section{Conclusion}
The foundation is solid and the path forward is clear. The hardest
cross-cutting concerns are already solved and covered by tests: secure
authentication, layered authorization, a reproducible schema, and an automated
quality pipeline. The matching algorithm that motivates the project is formally
specified and ready to build on the data structures already in place. The
remaining phases are well-scoped and ordered so each one builds on the last.

The full thesis will demonstrate that a configurable, multi-criteria matching
mechanism can be implemented securely and evaluated rigorously on real data. Once
the matching engine is complete and the evaluation is run, the EMS Portal will
show whether intelligent provider selection—optimized across skill, availability,
cost, location, and reputation—produces better outcomes than keyword search alone.
